\documentclass[./Main.tex]{subfiles}
\graphicspath{{./Fig/}}
\begin{document}

\chapter{解析手法}
\label{chap:simulation}

\section{解析条件}
\cref{tab:condition}に条件を示す。

\begin{table}[htbp]
    \centering
     \caption{実験条件}
    \begin{tabular}{cc}
    \hline
    項目     & 　内容 \\\hline
    ケース1     & 乱流 \\
    ケース2  &  層流 \\\hline
    \end{tabular}
       \label{tab:condition}
\end{table}

\section{はじめに}
陰解法を用いた場合，行列の反転，すなわち連立一次方程式を解く必要がある．
連立一次方程式の解法には直接法と反復法がある．直接法では厳密に解が求めら
れるが，計算量が多い．反復法は反復計算により近似解を求める方法であり，係
数行列の特性に応じて，高速に解を得られることがある．
CFDでは反復法が多く用いられる．行列の特性によって収束の速度に違いが出る
と考えられるので，各手法を適用し，最適なものを選択する．ここでは\cref{eq:linear}で
表される$n$次元の連立一次方程式を解くことを考える．
\begin{equation}
 A\bm{x}=\bm{b}
\label{eq:linear}
\end{equation}
$A$は$n\times n$の係数行列，$x$は$n\times1$の解ベクトル，$b$は
$n\times1$のベクトルを表す．

\cref{eq:linear}

この章では主に文献\cite{issa1986}を参考にした．

\section{直接法}
連立1次方程式を行列演算によって直接解く，直接法であるガウスの消去法につ
いて述べる．
\subsection{ガウスの消去法}
ガウスの消去法は．係数行列$A$を一度上三角行列に変形し，それから単位行列
$I$を求めることで解を得る．

ガウスの消去法による演算量は$O(n^3)$に比例して増大するため，次元の大きな
連立方程式を解くのは非現実的となる．このためCFDでは反復法が多く用いられ
る．

\section{反復法}
反復法では，初期値として適当な値を与え，真の解に収束する近似解を反復計算により求める．ただし，係数行列
によっては反復計算が収束せず，解を得ることができない場合がある．反復解が
収束するためには係数行列が対角優位であることが必要である．対角優位とは，
行列の対角成分が他の成分に比べて，その絶対値が大きいような行列である．
\begin{equation}
 |a_{ii}|>\sum_{i\neq j}|a_{ij}|, \;\;(i=1,...,n)
\end{equation}
ここで$a_{ij}$は係数行列の要素．

\section{乱流解析}
文献\cite{Kajishima}

\ifSubfilesClassLoaded{
\bibliographystyle{cs2srt}
\bibliography{./mybib.bib}
}{}

\end{document}


